\subsection{群胚的下降}

设 X 和 Y 是拓扑空间，设 \(f\colon\mathrm{X}\to\mathrm{Y}\) 是连续映射。设 \(p_{1}\) 和 \(p_{2}\) 是 \(\mathrm{X}\times_{\mathrm{Y}}\mathrm{X}\) 到 X 的两个投影，\(\operatorname{Coeg}(f)\) 是群胚态射对 \((\varpi(p_{1}),\varpi(p_{2}))\) 的余均衡群胚，这两个态射从 \(\varpi(\mathrm{X}\times_{\mathrm{Y}}\mathrm{X})\) 到 \(\varpi(\mathrm{X})\)，而 \(\gamma\colon\varpi(\mathrm{X})\to\operatorname{Coeg}(f)\) 是规范群胚态射（II, 第~199~页，定义 2）。由于 \(f\circ p_{1}=f\circ p_{2}\)，有 \(\varpi(f)\circ\varpi(p_{1})=\varpi(f)\circ\varpi(p_{2})\)。因此根据余均衡子的泛性质（II, 第~199~页，命题 3），存在唯一的群胚态射 \(\varpi'(f)\colon\operatorname{Coeg}(f)\to\varpi(\mathrm{Y})\)，使得 \(\varpi(f)=\varpi'(f)\circ\gamma\)。

Coeg(f) 的顶点集是集合 X = Vert( \(\varpi(X)\) ) 按 f 定义的等价关系所得的商集（II, 第~200~页，注 1）。它可识别为 f(X)。

命题 8. --- 设 X 和 Y 是非空拓扑空间，f 是从 X 到 Y 的连续映射。假设 X 局部道路连通，Y 连通，且 f 严格且满射。那么群胚 Coeg(f) 是传递的。

设 \(\Gamma\) 是这样的箭图：其顶点集为 X，箭集为 \(\mathrm{Fl}(\varpi(\mathrm{X}))\) 与 \(\mathrm{X}\times_{\mathrm{Y}}\mathrm{X}\) 的不交并；源映射和目标映射是 \(\varpi(\mathrm{X})\) 的源映射和目标映射（作用于 \(\mathrm{Fl}(\varpi(\mathrm{X}))\)），以及映射 \(p_1\) 和 \(p_2\)（作用于 \(\mathrm{X}\times_{\mathrm{Y}}\mathrm{X}\)）。根据对 \((\varpi(\mathrm{pr}_1), \varpi(\mathrm{pr}_2))\) 的框架的定义（II, 第~185~页，定义 3）以及 II, 第~200~页注 2，\(\operatorname{Coeg}(f)\) 的轨道集可识别为箭图 \(\Gamma\) 的连通分支集。由于空间 X 非空，只需证明图 \(\Gamma\) 是连通的。

\(\Gamma\) 的连通分支关于等价关系“存在道路连接 \(x\) 与 \(x'\)”是饱和的，因此由于 X 局部道路连通，它们在 X 中是开集。于是它们也是闭集。它们关于 f 定义的等价关系 R 也是饱和的。根据假设，映射 f 诱导出从 X/R 到 Y 的同胚，所以每个 \(\Gamma\) 的连通分支在 f 下的像都是 Y 的开闭子集；由于空间 Y 连通，这个像必为 Y。

设 C 是 \(\Gamma\) 的一个连通分支，设 \(x\) 是 X 中一点。根据前述内容，存在 \(x' \in \mathrm{C}\)，使得 \(f(x') = f(x)\)。根据箭图 \(\Gamma\) 的定义，有 \(x' \in \mathrm{C}\)。因此 \(\mathrm{C} = \mathrm{X}\)，箭图 \(\Gamma\) 从而是连通的。

命题 9. --- 设 X 和 Y 是拓扑空间，f 是从 X 到 Y 的连续映射。假设 X 局部道路连通，Y 是可解圈的，且映射 f 严格且满射。那么，群胚 \(\varpi(\mathrm{Y})\) 由 \(\varpi(\mathrm{X})\) 在 \(\varpi(f)\) 下的像生成。

不妨设空间 X 和 Y 非空。令 G 为 \(\varpi(\mathrm{Y})\) 中由 \(\varpi(\mathrm{X})\) 在 \(\varpi(f)\) 下的像生成的子群胚。由于映射 f 满射，G 的顶点集等于 Y。根据命题 8，\(\operatorname{Coeg}(f)\) 是传递群胚。由于 \(\varpi'(f)\) 在顶点上诱导恒等，\(\operatorname{Coeg}(f)\) 在 \(\varpi'(f)\) 下的像是 \(\varpi(\mathrm{Y})\) 的传递子群胚。\(\varpi(\mathrm{X})\) 在 \(\gamma\) 下的像生成 \(\operatorname{Coeg}(f)\)（II, 第~200~页，推论）；由于有 \(\varpi(f) = \varpi'(f) \circ \gamma\)，群胚 G 是传递的。

设 \(y_{0}\) 是 Y 中一点，令 H 表示子群 \(G_{y_{0}}\)，它是 \(\pi_{1}(Y, y_{0})\) 的子群。根据 IV, 第~342~页定理 1，存在 Y 的连通覆叠（T, p）以及一点 \(t_{0}\)，它属于纤维 \(T_{y_{0}}\)，且其固定子为 H。

若 \(x\) 是 X 中一点，则集合 \(\mathrm{Fl}_{y_0,f(x)}(\mathrm{G})\) 非空，因为 G 是传递的。对于 \(u\in \mathrm{Fl}_{y_0,f(x)}(\mathrm{G})\)，点 \(t_0\cdot u\) 是纤维 \(\mathrm{T}_{f(x)}\) 中的一点，并且与 \(u\) 无关，因为群 H（G 在 \(y_0\) 处的各向同性群）固定 \(t_0\)。将此点记为 \(\sigma (x)\)；将由此定义的映射 \(\sigma \colon \mathrm{X}\to \mathrm{T}\) 记为，并令 \(s\colon \mathrm{X}\to \mathrm{X}\times_{\mathrm{Y}}\mathrm{T}\) 为映射 \(x\mapsto (x,\sigma (x))\)。按构造，映射 s 与 \(\varpi (\mathrm{X})\) 在 X 和 \(\mathrm{X}\times_{\mathrm{Y}}\mathrm{T}\) 上的规范作用相容。由 III, 第~312~页引理 1，s 是连续的。因此映射 \(\sigma\) 连续；此外，由于 \(\sigma (x)\) 只依赖于 \(f(x)\)，它与 f 定义的等价关系相容。由于 f 严格且满射，存在唯一连续映射 \(\overline{\sigma}\colon \mathrm{Y}\to \mathrm{T}\)，使得 \(\overline{\sigma}\circ f = \sigma\)，所以覆叠 T 有截面。由于 T 连通，映射 \(p\colon \mathrm{T}\to \mathrm{Y}\) 是同胚（I, 第~31~页，命题 6 的推论 4），并且 \(\mathrm{H} = \pi_1(\mathrm{Y},y_0)\)，从而 \(\mathrm{G}_{y_0} = \pi_1(\mathrm{Y},y_0)\)。由于 G 传递，得到 \(\mathrm{G} = \varpi (\mathrm{Y})\)。

命题 10. --- 设 X 和 Y 是拓扑空间，f 是从 X 到 Y 的连续映射。假设空间 X 是可解圈的，空间 \(X\times_{Y}X\) 局部道路连通，并且关于 f 在 X 的覆叠上的每个下降数据都是有效的，商空间是 Y 的覆叠。那么群胚态射 \(\varpi'(f)\) 从 Coeg(f) 到 \(\varpi(Y)\)，是单射。

在命题的假设下，\(f\) 是严格的（IV, 第~394~页，注）；此外不妨设它是满射的。

引理 2. --- 保持命题的记号和假设。设 T 是一个集合，赋予它相对于映射 q: T → Y 的群胚 Coeg(f) 作用。那么 T 上存在唯一拓扑，使得 q 使 T 成为 Y 的覆叠，并且对于以 X 中 c 为起点的每个道路类以及 T 中每个点 t，都有 \(t \cdot \gamma(c) = t \cdot f_{*}(c)\)。特别地，Coeg(f) 在 T 上的作用可通过群胚 \(\varpi(Y)\) 的作用分解。

这种 T 上拓扑的唯一性由 III, 第~312~页引理 1 得出；证明其存在性。给集合 \(X\times_{Y}T\) 赋予 \(\varpi(\mathrm{X})\) 的作用，规定 \((x,t)\cdot u=(x\cdot u,t\cdot\gamma(u))\)，其中 \((x,t)\in\mathrm{X}\times_{\mathrm{Y}}\mathrm{T}\)，且 u 是 \(\varpi(\mathrm{X})\) 中源为 x 的每个箭。由于空间 X 是可解圈的，此作用没有局部单值性（III, 第~313~页，注）。因此，在集合 \(X\times_{Y}T\) 上存在唯一拓扑，使 X-空间 \((\mathrm{X}\times_{\mathrm{Y}}\mathrm{T},\mathrm{pr}_{1})\) 成为 X 的覆叠，并使 \(\varpi(\mathrm{X})\) 在此覆叠上的规范作用与给定作用重合（III, 第~313~页，命题 3）。记此覆叠为 \((Z,p)\)。证明映射 \(\tau: ((x_{1},t),x_{2})\mapsto(x_{2},t)\)，从 \(Z\times_{Y}X\) 到 Z，是关于映射 f 在 Z 上的下降数据。只需验证它连续，IV, 第~382~页定义 1 的其他条件显然成立。

映射 \(\tau\) 是映射 \(\mathrm{pr}_2\colon Z\times_Y X\to X\) 到覆叠 \(Z\)（其底空间为 \(X\)）的提升。由于空间 \(X\times_Y X\) 局部道路连通，作为它的覆叠的空间 \(Z\times_Y X\) 也局部道路连通。根据推论（III, 第~269~页），要证明映射 \(\tau\) 连续，只需证明它沿道路连续。

设 \(\widetilde{c} = ((c,g),c')\) 是 \(Z\times_{Y}X\) 中的一条道路，证明映射 \(\tau \circ \widetilde{c}\colon t\mapsto (c'(t),g(t))\) 从 I 到 Z 连续。对于每个 \(s\in [0,1]\)，记 \(c_s\) 和 \(c_s'\) 为 X 中分别由 \(t\mapsto c(st)\) 和 \(t\mapsto c'(st)\) 定义的道路。对于所有 \(s\in [0,1]\)，于是有 \(c(s) = c(0)\cdot [c_s]\)、\(c'(s) = c'(0)\cdot [c_s']\) 以及 \((c,g)(s) = (c(0),g(0))\cdot [c_s]\)，其中 \([u]\) 表示道路 u 的严格同伦类（III, 第~304~页，注）。映射 \(t\mapsto (c_s(t),c_s'(t))\) 是 \(X\times_{Y}X\) 中的一条道路；根据余均衡群胚 Coeg(f) 的定义，有 \(\gamma ([c_s]) = \gamma ([c_s'])\)，并且 \(g(s) = g(0)\cdot \gamma ([c_s]) = g(0)\cdot \gamma ([c_s'])\)。根据 \(\varpi (X)\) 在 Z 上的作用定义，因而有

\[
\begin{array}{r} (\tau \circ \widetilde {c}) (s) = (c ^ {\prime} (s), g (s)) = (c ^ {\prime} (0) \cdot [ c _ {s} ^ {\prime} ], g (0) \cdot \gamma ([ c _ {s} ^ {\prime} ])) \\ = (c ^ {\prime} (0), g (0)) \cdot [ c _ {s} ^ {\prime} ] = (\tau \circ \widetilde {c}) (0) \cdot [ c _ {s} ^ {\prime} ]. \end{array}
\]

这说明 \(\tau \circ \widetilde{c}\) 是道路 \(c'\) 到 Z 的连续提升（同上）。因此，映射 \(\tau\) 沿道路连续，因而连续；它是 X-空间 Z 关于 \(f\) 的下降数据。

记 \(R_{\tau}\) 为下降数据 \(\tau\) 定义的等价关系。由于映射 f 满射，映射 \(pr_{2}: Z \to T\) 通过取商诱导出从 \(Z/R_{\tau}\) 到 T 的双射。通过结构传递，将由 \(Z/R_{\tau}\) 的拓扑得到的拓扑赋予 T，使得 \((T,q)\) 成为 Y-空间。根据假设，T 因而是 Y 的覆叠，且图表

\[
\begin{array}{c} \mathrm{Z} \xrightarrow {\mathrm{pr} _ {2}} \mathrm{T} \\ p \Big \downarrow \qquad \qquad \qquad \qquad \Big \downarrow q \\ \mathrm{X} \xrightarrow {f} \mathrm{Y} \end{array}
\]

是笛卡尔方形。

设 \((x,t)\) 是 Z 中一点，设 c 是 X 中以 x 为起点的一条道路，并设 \(\widetilde{c}\) 是以 \((x,t)\) 为起点的 c 到 Z 的连续提升。道路 \(pr_{2} \circ \widetilde{c}\) 是道路 \(f \circ c\) 在 Y 中以 t 为起点的提升，因此 \(t \cdot \gamma([c]) = t \cdot [f \circ c] = t \cdot f_{*}([c])\)。引理得证。

现在证明命题。设 \(u\) 和 \(v\) 是 Coeg( \(f\) ) 的两个箭，它们在 \(\varpi(Y)\) 中的像相等。由于群胚态射 \(\varpi'(f)\) 在顶点集上是恒等的，箭 \(u\) 和 \(v\) 有相同的起点和终点。记 \(y\) 为 \(u\) 的起点；令 T 为 Coeg( \(f\) ) 中以 y 为源的箭集，并令 \(q: \mathrm{T} \to \mathrm{Y}\) 为目标映射在 T 上的限制。群胚 Coeg( \(f\) ) 关于 \(q\) 通过右合成作用于集合 T。根据引理 2，\(u\) 和 \(v\) 在 T 上的作用相同。因此有 \(u = e_y \cdot u = e_y \cdot v = v\)。这证明群胚态射 \(\varpi'(f)\) 是单射。

定理 1. --- 设 X 和 Y 是拓扑空间，设 \(f: X \to Y\) 是连续且满射的映射。假设空间 X 和 Y 都是可解圈的。最后，假设下列性质之一成立：

\begin{enumerate}
\def\labelenumi{(\roman{enumi})}
\item
  映射 f 逆紧、分离，具有局部连通的纤维，且空间 \(X\times_{Y}X\) 局部道路连通；
\item
  映射 f 逆紧、分离，具有有限纤维，对角线 \(\Delta_{X}\) 在 \(X \times_{Y} X\) 中是开集，且其补集是局部连通的；
\item
  映射 f 是开映射并具有道路提升性质。
\end{enumerate}

那么，群胚态射 \(\varpi'(f)\) 是从群胚 Coeg(f) 到庞加莱群胚 \(\varpi(Y)\) 的同构。

首先注意，在这些假设下，映射 \(f\) 是满射且严格的（I, 第~18~页，例 2）。此外，关于 \(f\) 在 X 的覆叠上的每个下降数据都是有效的，商空间是 Y 的覆叠；这在假设（i）和（ii）下由 IV, 第~393~页命题 4 的推论 2 得出，在假设（iii）下由 IV, 第~394~页命题 7 得出。根据命题 10，群胚态射 \(\varpi'(f)\) 因而是单射，且其像是 \(\varpi(Y)\) 的子群胚。

根据 IV, 第~395~页命题 9，在假设（i）和（ii）下，此像等于 \(\varpi(\mathrm{Y})\)；在假设（iii）下也如此，因为此时群胚态射 \(\varpi(f)\) 是满射的。

因此，\(\varpi'(f)\) 是同构。

例。--- 下面是两个满足定理 1 假设的例子。

\begin{enumerate}
\def\labelenumi{\arabic{enumi})}
\item
  设 Y 是可解圈拓扑空间。设 \((\mathrm{A}_{i})_{i\in\mathrm{I}}\) 是由闭集构成的 Y 的局部有限覆盖。假设对每个 \(i\in I\)，空间 \(A_{i}\) 是可解圈的，并且对每一对 \((i,j)\in\mathrm{I}\times\mathrm{I}\)，空间 \(A_{i}\cap A_{j}\) 局部道路连通。可以取 X 为族 \((\mathrm{A}_{i})_{i\in\mathrm{I}}\) 的和空间，取 \(f:X\to Y\) 为由规范嵌入族导出的映射。
\item
  设 G 是作用在可解圈拓扑空间 X 上的离散群，且作用是适当的；令 Y = X/G，并令 \(f: X \to Y\) 为规范映射。它是开映射（TG, III, 第~10~页，引理 2），并根据 III, 第~287~页定理 4 具有道路提升性质。根据 IV, 第~349~页命题 8 b)，空间 Y 是可解圈的。由从 G × X 到 \(X\times X\) 的映射 \((g, x) \mapsto (g \cdot x, x)\) 给出的映射是逆紧的，因而严格（I, 第~18~页，例 2），且其像是 \(X\times_{Y}X\)。于是根据 III, 第~261~页命题 8，\(X\times_{Y}X\) 局部道路连通。
\end{enumerate}

